\section{Introducción}

La Temperature Control Lab (TCLab) es una planta térmica didáctica montada sobre un Arduino
Leonardo. Dispone de dos transistores que actúan como calentadores y de dos sensores de temperatura
adheridos a ellos; la potencia de cada calentador se comanda como un porcentaje de su valor máximo
y la temperatura se mide en grados Celsius. En este trabajo se utilizó únicamente el primer par
calentador--sensor, de modo que la variable manipulada es $Q_1$, expresada en porcentaje, y la
variable controlada es $T_1$.

El problema de control consiste en llevar $T_1$ desde la temperatura ambiente hasta una referencia
de \SI{60}{\celsius} y mantenerla en ese valor, respetando el límite de potencia
$Q_1 \in [0, 90]\,\si{\percent}$ impuesto para proteger el equipo. Los requisitos de desempeño
establecidos fueron un tiempo de establecimiento $t_s \leq \SI{600}{\second}$ y un sobreimpulso
$M_p \leq \SI{5}{\percent}$.

El objetivo del trabajo no se limita a obtener un controlador que funcione en simulación, sino a
recorrer un flujo completo de diseño basado en modelos (\emph{Model-Based Design}, MBD). En este
enfoque el modelo ejecutable es el eje del desarrollo, y entre el diseño matemático y el sistema
físico se intercalan etapas sucesivas de verificación que comprueban aspectos distintos de la
solución \cite{mw_hil}. El flujo seguido fue

\begin{center}
modelado $\rightarrow$ identificación $\rightarrow$ diseño $\rightarrow$ MIL $\rightarrow$
código $\rightarrow$ SIL $\rightarrow$ PIL $\rightarrow$ implementación física.
\end{center}

Cada etapa acota el riesgo de la siguiente. La evaluación \emph{Model-in-the-Loop} (MIL) comprueba
el desempeño del algoritmo contra un modelo de la planta; \emph{Software-in-the-Loop} (SIL)
verifica que el código generado se comporte como el modelo del que proviene;
\emph{Processor-in-the-Loop} (PIL) verifica que ese código conserve su comportamiento al ejecutarse
sobre el procesador objetivo; y solo entonces el controlador se cierra sobre la planta real. El
documento sigue ese mismo orden.
